\documentclass[10pt]{letter}

\usepackage[margin=0.8in]{geometry}
\usepackage[hidelinks]{hyperref}
\usepackage[T1]{fontenc}
\usepackage{lmodern}

\signature{Soroush Vahidi\\
Ph.D. Candidate in Computer Science\\
Department of Computer Science\\
Ying Wu College of Computing\\
New Jersey Institute of Technology\\
\href{mailto:sv96@njit.edu}{sv96@njit.edu}}

\address{Soroush Vahidi\\
Department of Computer Science\\
Ying Wu College of Computing\\
New Jersey Institute of Technology\\
Newark, NJ 07102, USA}

\date{September 28, 2026}

\begin{document}

\begin{letter}{Editors of \emph{Performance Evaluation}}

\opening{Dear Editors of \emph{Performance Evaluation},}

I am pleased to submit the manuscript ``How Portable Are LLM-Serving
Scheduler Rankings Across Workloads, Operating Regions, and Metrics?'' for
consideration as an original research article in \emph{Performance
Evaluation}.

The manuscript studies whether comparative rankings of LLM-serving
schedulers remain portable when the workload source, operating region, or
evaluation metric changes. Its main methodological contribution is to treat
comparative scheduler-ranking portability itself as an object of performance
evaluation, rather than to propose another scheduling policy. The study uses
three independently released workload sources, six calibrated operating
regions, and 9,360 unique configurations, each executed twice.

The results are intentionally balanced. Cross-source rankings are stable for
some source pairs, but 36/990 pairwise comparisons show statistically
supported practical reversals, all concentrated in one policy pair with
BurstGPT on one side. Cross-metric portability is substantially lower:
median Kendall's tau-b is 0.419 and top-policy agreement is 31.9\%. A
selected physical validation did not reproduce the largest simulated
reversal, while the stable-control comparison preserved its qualitative
ordering.

The manuscript fits \emph{Performance Evaluation}'s scope through its focus
on performance evaluation of computing systems, scheduling and resource
allocation, trace-driven simulation, workload characterization, statistical
analysis, and AI-based services. Relative to workload-sensitive dispatching
studies, including the data-center dispatching analysis of Yildiz et al.\
(2026), LSSP evaluates the
portability of an entire comparative LLM-serving scheduler ranking across
independent workload sources, operating regions, and metrics, with explicit
ranking uncertainty, practical reversal detection, benchmark
sample-complexity analysis, and selected physical validation.

Artifacts are publicly available through GitHub, Hugging Face, and Zenodo:
\begin{quote}\small
\url{https://github.com/SoroushVahidi/llm-serving-scheduler-robustness-benchmark}\\
\url{https://huggingface.co/datasets/SoroushVahidi/llm-serving-scheduler-portability}\\
\url{https://doi.org/10.5281/zenodo.22306798}
\end{quote}
As described in the manuscript, the release includes source code, provenance,
selected canonical result tables, and the reduced RQ6 analysis result; it
does not claim to publish all raw per-execution records.

This manuscript is original, is not under consideration elsewhere, and I
approve its submission as the sole author. Relevant competing-interest
information and other required declarations are included in the manuscript.

\closing{Sincerely,}

\end{letter}

\end{document}
